% Generated from lessons/0011-2-multivariable-continuous-functions.html; edit the HTML source, then rerun the converter.
\section{多元连续函数}
\lessonmark{多元连续函数}

这一节把一元函数的极限、连续、四则运算和复合函数推广到 \(R^n\)。期末最常考的是：多元极限存在性、路径反例、累次极限、连续性的判定与向量值函数的分量判别。

\subsection*{本节地图}

\begin{conceptbox}
\textbf{多元函数}
输入从一个数变成一个点 \(x=(x_1,\ldots,x_n)\)，输出仍是一个实数。

\[
          f:D\subset R^n\to R,\qquad x\mapsto z=f(x)
        \]
\end{conceptbox}

\begin{conceptbox}
\textbf{多元极限}
要求 \(x\) 以任意方式趋于 \(x_0\) 时，函数值都趋于同一个数。
\end{conceptbox}

\begin{conceptbox}
\textbf{连续与映射}
连续就是 \(\lim_{x\to x_0}f(x)=f(x_0)\)。向量值函数可按坐标分量判定。
\end{conceptbox}

教材对应：《数学分析》第三版下册第十一章 §2“多元连续函数”。辅助参考：\href{https://www.math.ucdavis.edu/~hunter/intro_analysis_pdf/ch13.pdf}{UC Davis Chapter 13}。

\subsection*{一、多元函数是什么}

一元函数是一个自变量控制一个函数值；多元函数是多个自变量共同控制函数值。例如理想气体中，压强 \(p\) 可由体积 \(V\) 和温度 \(T\) 决定；圆台体积可由两个底半径和高度决定。

\[
      f:D\subset R^n\to R,\qquad
      x=(x_1,\ldots,x_n)\mapsto z=f(x)
    \]

\begin{center}
\small
\begin{tabularx}{\linewidth}{|>{\raggedright\arraybackslash}p{0.18\linewidth}|>{\raggedright\arraybackslash}p{0.42\linewidth}|>{\raggedright\arraybackslash}X|}
\hline
\textbf{术语} & \textbf{意思} & \textbf{考试关注点} \\ \hline
定义域 \(D\) & 允许输入的点 \(x\) 的全体。 & 根号、分母、对数、反三角函数会限制定义域。 \\ \hline
值域 \(f(D)\) & 所有可能的函数值。 & 本节通常不重点求值域。 \\ \hline
图像 & \(\Gamma=\{(x,z)\in R^{n+1}:z=f(x),x\in D\}\)。 & 二元函数图像通常是 \(R^3\) 中曲面。 \\ \hline
\end{tabularx}
\end{center}

\subsection*{二、多元极限：定义与判题}

设 \(D\subset R^n\)，\(x_0\) 是 \(D\) 的聚点，函数 \(f\) 在 \(D\setminus\{x_0\}\) 上有定义。若存在实数 \(A\)，使得：

\[
      \forall \varepsilon>0,\ \exists \delta>0,\quad
      0< |x-x_0|< \delta,\ x\in D
      \Longrightarrow |f(x)-A|<\varepsilon,
    \]

则称 \(f(x)\) 当 \(x\to x_0\) 时极限为 \(A\)，记作：

\[
      \lim_{x\to x_0} f(x)=A.
    \]

\begin{notebox}
多元极限比一元极限更严：不是只从左、右靠近，而是从平面或空间中的所有路径靠近。只要有两条路径给出不同极限，二重极限就不存在。
\end{notebox}

\subsubsection*{判断二重极限的考试流程}

\begin{enumerate}
 \item 先看目标点处能否直接用连续性代入，例如多项式、分母非零的有理式、连续函数的复合。
 \item 若不能直接代入，先用轴线、直线或累次极限猜可能的极限值 \(A\)。
 \item 怀疑不存在时，找两条路径给出不同极限；直线不够时试曲线。
 \item 怀疑存在时，必须回到估计或夹逼，把 \(|f(x,y)-A|\) 压到 \(0\)。
 \end{enumerate}

\subsubsection*{证明极限存在的常用方法}

\begin{center}
\small
\begin{tabularx}{\linewidth}{|>{\raggedright\arraybackslash}p{0.18\linewidth}|>{\raggedright\arraybackslash}p{0.42\linewidth}|>{\raggedright\arraybackslash}X|}
\hline
\textbf{方法} & \textbf{典型写法} & \textbf{适用信号} \\ \hline
连续性法 & 若 \(f\) 在 \((x_0,y_0)\) 连续，则 \(\lim f=f(x_0,y_0)\)。 & 多项式、三角函数、指数对数、分母不为 \(0\) 的商、连续函数复合。 \\ \hline
半径估计法 & 令 \(r=\sqrt{(x-x_0)^2+(y-y_0)^2}\)，把 \(|f(x,y)-A|\) 控制成 \(Cr^\alpha\)。 & 分子分母含 \(x^2+y^2\)、高阶小量、需要证明极限为 \(0\)。 \\ \hline
夹逼法 & 找到 \(0\le |f(x,y)-A|\le h(x,y)\)，且 \(h(x,y)\to0\)。 & 含 \(\sin,\cos\) 或振荡项，常用 \(|\sin t|\le1\)、\(|\cos t|\le1\)。 \\ \hline
极坐标法 & 在原点令 \(x=r\cos\theta,\ y=r\sin\theta\)。若 \(f=A+r^\alpha B(\theta)\)，且 \(B(\theta)\) 有界，则极限为 \(A\)。 & 原点处的齐次式、含 \(x^2+y^2\) 的分母、想看角度 \(\theta\) 是否影响结果。 \\ \hline
运算法则 & 若 \(f\to A,\ g\to B\)，则 \(f+g\to A+B,\ fg\to AB,\ f/g\to A/B\)，其中 \(B\ne0\)。 & 函数能拆成几个已知极限的和、积、商。 \\ \hline
\end{tabularx}
\end{center}

\[
      |f(x,y)-A|\le C\left(\sqrt{(x-x_0)^2+(y-y_0)^2}\right)^\alpha,\quad \alpha>0
      \Longrightarrow
      \lim_{(x,y)\to(x_0,y_0)}f(x,y)=A.
    \]

\begin{notebox}
证明存在不能只看几条路径。路径只覆盖趋近方式的一小部分；真正的存在证明要能管住所有趋近方式。
\end{notebox}

\subsubsection*{证明极限不存在的常用方法}

最标准的动作是找两条趋近路径，让函数值趋于不同结果。例如：

\[
      f(x,y)=\frac{xy}{x^2+y^2}.
    \]

沿 \(y=0\) 趋于 \((0,0)\)，有 \(f(x,0)=0\)；沿 \(y=x\) 趋于 \((0,0)\)，有 \(f(x,x)=1/2\)。两条路径不同，所以极限不存在。

\begin{center}
\small
\begin{tabularx}{\linewidth}{|>{\raggedright\arraybackslash}p{0.18\linewidth}|>{\raggedright\arraybackslash}p{0.42\linewidth}|>{\raggedright\arraybackslash}X|}
\hline
\textbf{反例路径} & \textbf{什么时候试} & \textbf{判定方式} \\ \hline
轴线 \(x=x_0\)、\(y=y_0\) & 先手检查，计算最简单。 & 两条轴线极限不同，二重极限不存在。 \\ \hline
直线 \(y-y_0=k(x-x_0)\) & 分子分母同阶，或出现 \(xy,x^2,y^2\)。 & 结果依赖 \(k\)，二重极限不存在。 \\ \hline
曲线 \(y-y_0=(x-x_0)^2\) 等 & 所有直线路径看起来相同，但式子里有 \(y-x^2\)、\(y^2-x\) 等结构。 & 曲线极限与直线极限不同，二重极限不存在。 \\ \hline
累次极限 & 需要快速排除存在性。 & 两个累次极限存在但不相等，则二重极限不存在。 \\ \hline
\end{tabularx}
\end{center}

\begin{notebox}
两个累次极限相等不能推出二重极限存在；但两个累次极限不相等，可以立刻推出二重极限不存在。
\end{notebox}

\subsection*{三、累次极限}

二重极限是 \((x,y)\) 一起趋近；累次极限是先让一个变量趋近，再让另一个变量趋近。

\[
      \lim_{x\to x_0}\lim_{y\to y_0}f(x,y),
      \qquad
      \lim_{y\to y_0}\lim_{x\to x_0}f(x,y).
    \]

\begin{center}
\small
\begin{tabularx}{\linewidth}{|>{\raggedright\arraybackslash}p{0.18\linewidth}|>{\raggedright\arraybackslash}p{0.42\linewidth}|>{\raggedright\arraybackslash}X|}
\hline
\textbf{现象} & \textbf{结论} & \textbf{考试提醒} \\ \hline
二重极限存在，且某一方向的内层极限在外层变量的去心邻域内存在 & 相应累次极限存在且等于二重极限。 & 这是教材定理 11.2.1 的核心。 \\ \hline
两个累次极限都存在且相等 & 不能推出二重极限存在。 & 累次极限只看两种特殊趋近顺序。 \\ \hline
两个累次极限存在但不相等 & 二重极限一定不存在。 & 这是很快的反证工具。 \\ \hline
\end{tabularx}
\end{center}

\subsubsection*{定理 11.2.1 的内容和证明思路}

若

\[
      \lim_{(x,y)\to(x_0,y_0)}f(x,y)=A,
    \]

且当 \(x\) 在 \(x_0\) 附近、\(x\ne x_0\) 时，内层极限

\[
      \varphi(x)=\lim_{y\to y_0}f(x,y)
    \]

存在，那么

\[
      \lim_{x\to x_0}\varphi(x)=A.
    \]

证明思路：二重极限给出 \(|f(x,y)-A|<\varepsilon/2\) 的小邻域。固定一个 \(x\) 后令 \(y\to y_0\)，不等式传给 \(\varphi(x)\)，于是 \(|\varphi(x)-A|\le\varepsilon/2\)。

\subsection*{四、连续性与运算}

若 \(f\) 在 \(x_0\) 有定义，且

\[
      \lim_{x\to x_0}f(x)=f(x_0),
    \]

则称 \(f\) 在 \(x_0\) 连续。用 \(\varepsilon-\delta\) 语言就是：

\[
      \forall\varepsilon>0,\ \exists\delta>0,\quad
      |x-x_0|<\delta
      \Longrightarrow |f(x)-f(x_0)|<\varepsilon.
    \]

一元连续函数的基本运算性质可以推广到多元函数：

\begin{center}
\small
\begin{tabularx}{\linewidth}{|>{\raggedright\arraybackslash}p{0.18\linewidth}|>{\raggedright\arraybackslash}p{0.42\linewidth}|>{\raggedright\arraybackslash}X|}
\hline
\textbf{性质} & \textbf{内容} & \textbf{证明思路} \\ \hline
和差积商 & 连续函数的和、差、积仍连续；分母非零时商连续。 & 由极限四则运算法则推出。 \\ \hline
复合 & 若 \(g:D\to\Omega\) 连续，\(f:\Omega\to R^m\) 连续，则 \(f\circ g\) 连续。 & 先用 \(f\) 的连续性找 \(\eta\)，再用 \(g\) 的连续性找 \(\delta\)。 \\ \hline
常见初等函数 & 多项式、根式、对数、三角函数在自然定义域内连续。 & 看作一元连续函数与多元连续函数的复合。 \\ \hline
\end{tabularx}
\end{center}

\subsection*{五、向量值函数}

向量值函数的输出不是一个实数，而是一个向量：

\[
      f:D\subset R^n\to R^m,\qquad
      f(x)=(f_1(x),f_2(x),\ldots,f_m(x)).
    \]

极限和连续都可按分量判断：

\[
      f(x)\to A=(A_1,\ldots,A_m)
      \Longleftrightarrow
      f_j(x)\to A_j,\quad j=1,\ldots,m.
    \]

\[
      f \text{ 在 } x_0 \text{ 连续}
      \Longleftrightarrow
      f_1,\ldots,f_m \text{ 都在 } x_0 \text{ 连续}.
    \]

证明思路：一边用 \(|f_j(x)-f_j(x_0)|\le |f(x)-f(x_0)|\)；另一边用

\[
      |f(x)-f(x_0)|
      \le\sum_{j=1}^{m}|f_j(x)-f_j(x_0)|.
    \]

\subsection*{六、即时判断}

\begin{quickcheck}
所有直线路径的极限都相同，是否一定能推出二重极限存在？

\tcblower
\textbf{答案：}不能
\end{quickcheck}

\begin{quickcheck}
若二重极限存在，且两个累次极限也存在，那么这三个极限的关系是什么？

\tcblower
\textbf{答案：}相等
\end{quickcheck}

\begin{quickcheck}
判断 \(f:D\subset R^n\to R^m\) 连续，最常用的拆法是什么？

\tcblower
\textbf{答案：}分量
\end{quickcheck}

\subsection*{七、课后题训练}

以下是第十一章 §2 课后题的完整训练版，表述做了复习化整理，公式和条件尽量保留。每题后面给出答案方向，便于不翻原教材也能复习。

\begin{exercisebox}
\textbf{习题 1\badge{自然定义域}}

确定下列函数的自然定义域：

\[
      \begin{aligned}
      &(1)\quad u=\ln(y-x)+\frac{x}{\sqrt{1-x^2-y^2}},\\
      &(2)\quad u=\frac1{\sqrt x}+\frac1{\sqrt y}+\frac1{\sqrt z},\\
      &(3)\quad u=\sqrt{R^2-x^2-y^2-z^2}+\sqrt{x^2+y^2+z^2-r^2}\quad(R> r),\\
      &(4)\quad u=\arcsin\frac{z}{x^2+y^2}.
      \end{aligned}
      \]

\begin{hintbox}
分别检查对数真数、分母、根号非负、反正弦自变量范围。

\[
        \begin{aligned}
        &(1)\ y> x,\quad x^2+y^2<1;\\
        &(2)\ x>0,\ y>0,\ z>0;\\
        &(3)\ r^2\le x^2+y^2+z^2\le R^2;\\
        &(4)\ x^2+y^2\ne0,\quad |z|\le x^2+y^2.
        \end{aligned}
        \]
\end{hintbox}
\end{exercisebox}

\begin{exercisebox}
\textbf{习题 2\badge{齐次结构}}

设 \(x>0\)，且

\[
        f\!\left(\frac yx\right)=\frac{x^3}{(x^2+y^2)^{3/2}}.
      \]

求一元函数 \(f(t)\)。

\begin{hintbox}
令 \(t=y/x\)，则 \(x^2+y^2=x^2(1+t^2)\)。

\[
          f(t)=\frac{1}{(1+t^2)^{3/2}}.
        \]
\end{hintbox}
\end{exercisebox}

\begin{exercisebox}
\textbf{习题 3\badge{代换求函数}}

若

\[
        z(x,y)=\sqrt y+f(\sqrt x-1),
      \]

且当 \(y=4\) 时 \(z=x+1\)。求 \(f(t)\) 和 \(z(x,y)\)。

\begin{hintbox}
令 \(t=\sqrt x-1\)，则 \(x=(t+1)^2\)。由 \(2+f(t)=x+1\)，得到：

\[
          f(t)=t^2+2t,\qquad z(x,y)=x+\sqrt y-1.
        \]
\end{hintbox}
\end{exercisebox}

\begin{exercisebox}
\textbf{习题 4\badge{原点极限存在性}}

讨论下列函数当 \((x,y)\to(0,0)\) 时的极限是否存在：

\[
      \begin{aligned}
      &(1)\ f(x,y)=\frac{x-y}{x+y},\\
      &(2)\ f(x,y)=\frac{xy}{x^2+y^2},\\
      &(3)\ f(x,y)=
      \begin{cases}
      1,&0< y< x^2,\\
      0,&\text{其他点},
      \end{cases}\\
      &(4)\ f(x,y)=\frac{x^3y^3}{x^4+y^4}.
      \end{aligned}
      \]

\begin{hintbox}
(1) 沿 \(y=0\) 与 \(x=0\) 得到不同结果，不存在。

(2) 沿 \(y=0\) 得 \(0\)，沿 \(y=x\) 得 \(1/2\)，不存在。

(3) 沿 \(y=x^2/2\) 得 \(1\)，沿 \(y=0\) 得 \(0\)，不存在。

(4) 存在且等于 \(0\)。可用 \(2x^2y^2\le x^4+y^4\) 得 \(|f(x,y)|\le |xy|/2\to0\)。
\end{hintbox}
\end{exercisebox}

\begin{exercisebox}
\textbf{习题 5\badge{极限性质证明}}

证明多元函数极限具有唯一性、局部有界性、局部保号性和局部夹逼性。

\begin{hintbox}
这些性质与一元函数完全同型。唯一性用 \(|A-B|\le |A-f(x)|+|f(x)-B|\)。局部有界性取 \(\varepsilon=1\)。局部保号性取 \(\varepsilon=|A|/2\)。夹逼性用不等式在去心邻域内传递并令极限趋近。
\end{hintbox}
\end{exercisebox}

\begin{exercisebox}
\textbf{习题 6\badge{四则运算法则}}

设 \(x\to x_0\) 时 \(f(x)\) 与 \(g(x)\) 的极限存在。证明：

\[
      \begin{aligned}
      &(1)\ \lim_{x\to x_0}(f(x)\pm g(x))=\lim_{x\to x_0}f(x)\pm\lim_{x\to x_0}g(x),\\
      &(2)\ \lim_{x\to x_0}(f(x)g(x))=\left(\lim_{x\to x_0}f(x)\right)\left(\lim_{x\to x_0}g(x)\right),\\
      &(3)\ \lim_{x\to x_0}\frac{f(x)}{g(x)}
      =\frac{\lim_{x\to x_0}f(x)}{\lim_{x\to x_0}g(x)},\quad \lim_{x\to x_0}g(x)\ne0.
      \end{aligned}
      \]

\begin{hintbox}
和差用三角不等式；乘法用拆项 \(fg-AB=f(g-B)+B(f-A)\)，并用局部有界性；商法先由保号性保证分母在邻域内不为 \(0\)，再转化为乘法。
\end{hintbox}
\end{exercisebox}

\begin{exercisebox}
\textbf{习题 7\badge{计算极限}}

求下列极限：

\[
      \begin{aligned}
      &(1)\ \lim_{(x,y)\to(0,1)}\frac{1-xy}{1+x^2+y^2},&
      &(2)\ \lim_{(x,y)\to(0,0)}\frac{1+x^2+y^2}{x^2+y^2},\\
      &(3)\ \lim_{(x,y)\to(0,0)}\frac{\sqrt{1+xy}-1}{xy},&
      &(4)\ \lim_{(x,y)\to(0,0)}\frac{x^2+y^2}{\sqrt{1+x^2+y^2}-1},\\
      &(5)\ \lim_{(x,y)\to(0,0)}\frac{\ln(x^2+e^{y^2})}{x^2+y^2},&
      &(6)\ \lim_{(x,y)\to(0,0)}\frac{\sin(x^3+y^3)}{x^2+y^2},\\
      &(7)\ \lim_{(x,y)\to(0,0)}\frac{1-\cos(x^2+y^2)}{(x^2+y^2)x^2y^2},&
      &(8)\ \lim_{\substack{x\to+\infty\\ y\to+\infty}}(x^2+y^2)e^{-(x+y)}.
      \end{aligned}
      \]

\begin{hintbox}
\[
        (1)\ \frac12,\quad
        (2)\ \text{发散到 }+\infty,\quad
        (3)\ \frac12,\quad
        (4)\ 2,\quad
        (5)\ 1,\quad
        (6)\ 0,\quad
        (7)\ \text{无有限极限},\quad
        (8)\ 0.
        \]

(3)(4) 用有理化；(5) 用 \(\ln(1+t)\sim t\) 与 \(e^{y^2}-1\sim y^2\)；(7) 沿 \(x=y=t\) 会发散。
\end{hintbox}
\end{exercisebox}

\begin{exercisebox}
\textbf{习题 8\badge{二重极限与累次极限}}

讨论下列函数在原点的二重极限和两个累次极限：

\[
      \begin{aligned}
      &(1)\ f(x,y)=\frac{x^2y^2}{x^2y^2+(x-y)^2},\\
      &(2)\ f(x,y)=\frac{x^2(1+x^2)-y^2(1+y^2)}{x^2+y^2},\\
      &(3)\ f(x,y)=x\sin\frac1y+y\sin\frac1x.
      \end{aligned}
      \]

\begin{hintbox}
(1) 二重极限不存在：沿 \(y=x\) 得 \(1\)，沿 \(y=0\) 得 \(0\)。两个累次极限都为 \(0\)。

(2) 二重极限不存在：沿 \(y=0\) 与 \(x=0\) 分别趋于 \(1\) 和 \(-1\)。两个累次极限分别为 \(1\) 与 \(-1\)。

(3) 二重极限存在且为 \(0\)，因为 \(|f(x,y)|\le |x|+|y|\)。但两个累次极限都不存在，因为内层会出现 \(\sin(1/y)\) 或 \(\sin(1/x)\) 的振荡。
\end{hintbox}
\end{exercisebox}

\begin{exercisebox}
\textbf{习题 9\badge{构造不连续点}}

验证下面函数在原点不连续，而在其他点连续：

\[
      f(x,y)=
      \begin{cases}
      \dfrac{2}{x^2}\left(y-\dfrac12x^2\right),&x>0,\ \dfrac12x^2< y\le x^2,\\[6pt]
      \dfrac{1}{x^2}(2x^2-y),&x>0,\ x^2< y<2x^2,\\[6pt]
      0,&\text{其他点}.
      \end{cases}
      \]

\begin{hintbox}
在原点沿抛物线带内可得到非零值，例如取 \(y=x^2\) 时函数值为 \(1\)，而沿其他许多路径函数值为 \(0\)，故原点不连续。其他点离分段边界有正距离，或落在分段边界上时左右表达式相接为同一值，可逐段验证连续。
\end{hintbox}
\end{exercisebox}

\begin{exercisebox}
\textbf{习题 10\badge{连续范围}}

讨论函数的连续范围：

\[
      f(x,y)=
      \begin{cases}
      \dfrac{x^2y}{x^2+y^2},&x^2+y^2\ne0,\\
      0,&x^2+y^2=0.
      \end{cases}
      \]

\begin{hintbox}
除原点外由四则运算连续。在原点有

\[
          |f(x,y)|=\frac{x^2|y|}{x^2+y^2}\le |y|\to0,
        \]

所以原点也连续。连续范围为整个 \(R^2\)。
\end{hintbox}
\end{exercisebox}

\begin{exercisebox}
\textbf{习题 11\badge{差商极限}}

设 \(f(t)\) 在区间 \((a,b)\) 上具有连续导数，\(D=(a,b)\times(a,b)\)。定义

\[
      F(x,y)=
      \begin{cases}
      \dfrac{f(x)-f(y)}{x-y},&x\ne y,\\
      f'(x),&x=y.
      \end{cases}
      \]

证明对任意 \(c\in(a,b)\)，有

\[
        \lim_{(x,y)\to(c,c)}F(x,y)=f'(c).
      \]

\begin{hintbox}
当 \(x\ne y\) 时，由 Lagrange 中值定理，存在 \(\xi\) 介于 \(x,y\) 之间，使 \(F(x,y)=f'(\xi)\)。当 \((x,y)\to(c,c)\) 时 \(\xi\to c\)，再用 \(f'\) 的连续性。
\end{hintbox}
\end{exercisebox}

\begin{exercisebox}
\textbf{习题 12\badge{分别连续 + Lipschitz}}

设二元函数 \(f(x,y)\) 在开集 \(D\subset R^2\) 内对变量 \(x\) 是连续的，并且对变量 \(y\) 满足 Lipschitz 条件：对任意 \((x,y'),(x,y'')\in D\)，有

\[
        |f(x,y')-f(x,y'')|\le L|y'-y''|.
      \]

证明 \(f(x,y)\) 在 \(D\) 上连续。

\begin{hintbox}
在点 \((x_0,y_0)\) 附近拆成两段：

\[
          |f(x,y)-f(x_0,y_0)|
          \le |f(x,y)-f(x,y_0)|+|f(x,y_0)-f(x_0,y_0)|.
        \]

第一项由 Lipschitz 条件控制，第二项由关于 \(x\) 的连续性控制。
\end{hintbox}
\end{exercisebox}

\begin{exercisebox}
\textbf{习题 13\badge{映射运算连续}}

证明：若 \(f,g:D\to R^m\) 都是连续映射，则映射 \(f+g\) 与实值函数 \(\langle f,g\rangle\) 在 \(D\) 上都是连续的。

\begin{hintbox}
按分量证明最省事。若 \(f=(f_1,\ldots,f_m)\)，\(g=(g_1,\ldots,g_m)\)，则 \(f+g=(f_1+g_1,\ldots,f_m+g_m)\)，每个分量连续。内积函数为

\[
          \langle f,g\rangle=\sum_{j=1}^{m}f_jg_j,
        \]

由实值连续函数的乘法与有限求和保持连续可得。
\end{hintbox}
\end{exercisebox}

\begin{exercisebox}
\textbf{习题 14\badge{复合映射连续}}

证明复合映射连续性定理：设 \(\Omega\subset R^k\)、\(D\subset R^n\) 为开集，\(g:D\to\Omega\)、\(f:\Omega\to R^m\)，且 \(g\) 在 \(D\) 上连续，\(f\) 在 \(\Omega\) 上连续，则

\[
        f\circ g:D\to R^m,\qquad x\mapsto f(g(x))
      \]

在 \(D\) 上连续。

\begin{hintbox}
任取 \(x_0\in D\)。由 \(f\) 在 \(g(x_0)\) 连续，给定 \(\varepsilon\) 可找到 \(\eta\)，使 \(|u-g(x_0)|<\eta\Rightarrow |f(u)-f(g(x_0))|<\varepsilon\)。再由 \(g\) 在 \(x_0\) 连续，找到 \(\delta\)，使 \(|x-x_0|<\delta\Rightarrow |g(x)-g(x_0)|<\eta\)。合起来即可。
\end{hintbox}
\end{exercisebox}

\subsection*{八、期末写法模板}

\begin{conceptbox}
\textbf{证极限存在}
把 \(|f(x)-A|\) 估计到关于 \(r=|x-x_0|\) 的表达式，并令 \(r\to0\)。
\end{conceptbox}

\begin{conceptbox}
\textbf{证极限不存在}
找两条路径趋近同一点，算出两个不同极限。
\end{conceptbox}

\begin{conceptbox}
\textbf{讨论累次极限}
先固定一个变量求内层极限，再对剩下变量求外层极限；不要把它自动当成二重极限。
\end{conceptbox}

\begin{conceptbox}
\textbf{证连续}
优先用连续函数运算、复合函数连续、向量值函数分量判别；原点特殊点再用定义或估计。
\end{conceptbox}

上一节：第十一章 §1：Euclid 空间上的基本定理。路线图：教材路线图。
